% \subsection{RMSNorm could induce TLCM}
One natural cause of TLCM is RMSNorm. RMSNorm normalizes the input to the attention and MLP sublayers in nearly all of the open-source models analyzed.  
%It operates by normalizing the residual stream to a fixed norm of $\sqrt{d_\text{m}}$, followed by scaling each component of the resulting vector with a learnable parameter. 
Formally, for an input $\mathbf{x} \in \mathbb{R}^{d_\text{m}}$ to an attention or MLP sublayers and a learnable parameter vector $\mathbf{g}$, RMSNorm is defined as:
$$\overline{\mathbf{x}} = \frac{\mathbf{x}}{\operatorname{RMS}(\mathbf{x})} \odot \mathbf{g} \qquad \qquad \operatorname{RMS}(\mathbf{x}) = \frac{1}{\sqrt{d_\text{m}}} \|\mathbf{x}\|_2,$$
where $d_\text{m}$ is the dimension of the residual stream. 

Crucially, the use of RMSNorm implies that both the attention and MLP sublayers have \emph{layernorm blindness}; they are blind to the norm of the residual stream. This blindness is significant because these sublayers predict contributions to the residual stream (i.e. $\mathbf{x}_i + \operatorname{Sublayer}(\mathbf{x}_i)$), whose relative impacts depend on the the residual streams current magnitude. 
%To maintain proportionality in their contributions when the residual stream norm is larger than expected, the contributions themselves must scale accordingly. 
Without visibility into the residual stream norm, attention and MLP sublayers risk under-contributing when the norm is high, which potentially leads to their contributions being overshadowed. This could encourage over-contribution behaviors followed by correction, which is consistent with TLCM. 
%This insight may interest practitioners focusing on training dynamics or model architecture, as addressing this inefficiency could provide opportunities for optimization.


However, we find that \emph{LayerNorm blindness alone does not fully explain the  TLCM}. This is because as discussed in Sec. \ref{sec:jacobian-selective}, the correction mechanism does not entirely reverse $\mathbf{a}_t$, contrary to what would be expected if RMSNorm were the primary cause. But perhaps more critically, models trained with alternative RMSNorm implementations continue to exhibit the mechanism:

\textbf{Gemma.} The sublayers in Gemma 1 were trained using pre-LayerNorm as described above. In contrast, Gemma 2 utilized a hybrid approach, employing both pre-LayerNorm and post-LayerNorm:
$$\mathbf{x}_{i+1} = \mathbf{x}_i + \operatorname{LayerNorm}(\operatorname{Sublayer}(\operatorname{LayerNorm}(\mathbf{x}_i))).$$
The addition of post-LayerNorm should, in principle, make the residual stream norm more predictable. However, empirical results show that the correction mechanism remains robust in Gemma 2. Refer to Figure \ref{fig:rmsnorm} for a comparison.

\textbf{OLMo.} The original OLMo models employed pre-LayerNorm exclusively. In the OLMo 2 series, pre-LayerNorm was replaced with post-LayerNorm and QK norm. Despite this architectural change, the correction mechanism persists strongly in OLMo 2, as shown in Figure \ref{fig:rmsnorm}.

% There are a variety of possible explanations for this phenomenon. For example, we would observe this phenomena if transformer blocks were communicating with each other: layer $i$ could write to some subspace while layer $i+1$ reads and deletes the same subspace.

